% Minimal LaTeX article scaffold for mathematical papers.
% Loads amsmath/amsthm/amssymb, sets up theorem environments with shared
% counter, and includes a basic title/abstract/MSC/keywords block.
%
% Compile with: pdflatex -> bibtex -> pdflatex -> pdflatex

\documentclass[11pt,a4paper]{article}

\usepackage[utf8]{inputenc}
\usepackage[T1]{fontenc}
\usepackage{lmodern}
\usepackage{microtype}

\usepackage{amsmath,amsthm,amssymb,mathtools}
\usepackage{bm}
\usepackage{hyperref}
\usepackage[capitalize,noabbrev]{cleveref}

% --- Theorem environments (shared counter) -----------------------------------
\theoremstyle{plain}
\newtheorem{theorem}{Theorem}[section]
\newtheorem{lemma}[theorem]{Lemma}
\newtheorem{proposition}[theorem]{Proposition}
\newtheorem{corollary}[theorem]{Corollary}

\theoremstyle{definition}
\newtheorem{definition}[theorem]{Definition}
\newtheorem{example}[theorem]{Example}
\newtheorem{assumption}[theorem]{Assumption}

\theoremstyle{remark}
\newtheorem{remark}[theorem]{Remark}

% --- Equation numbering: per-section ----------------------------------------
\numberwithin{equation}{section}

% --- Operators ---------------------------------------------------------------
\DeclareMathOperator*{\argmin}{arg\,min}
\DeclareMathOperator*{\argmax}{arg\,max}
\DeclareMathOperator{\supp}{supp}
\DeclareMathOperator{\Var}{Var}
\DeclareMathOperator{\Cov}{Cov}

% --- Frequently-used sets ----------------------------------------------------
\newcommand{\R}{\mathbb{R}}
\newcommand{\N}{\mathbb{N}}
\newcommand{\Z}{\mathbb{Z}}
\newcommand{\Q}{\mathbb{Q}}
\newcommand{\C}{\mathbb{C}}
\newcommand{\E}{\mathbb{E}}
\newcommand{\Prb}{\mathbb{P}}

% --- Convenience -------------------------------------------------------------
\newcommand{\norm}[1]{\left\lVert #1 \right\rVert}
\newcommand{\abs}[1]{\left\lvert #1 \right\rvert}
\newcommand{\inner}[2]{\left\langle #1, #2 \right\rangle}

% =============================================================================
\title{Paper Title}
\author{Author Name\thanks{Affiliation, email@example.com}}
\date{\today}

\begin{document}
\maketitle

\begin{abstract}
Abstract goes here. Keep it under 200 words. State the problem, the approach,
the main result, and the contribution.
\end{abstract}

\noindent\textbf{Keywords:} keyword1, keyword2, keyword3.

\noindent\textbf{MSC2020:} 00A00, 11B11, 22C22.

% =============================================================================
\section{Introduction}\label{sec:intro}

Motivation, prior work, contributions.

\paragraph{Notation.}
Throughout, $\R$ denotes the real numbers and $\norm{\cdot}$ the Euclidean
norm. Vectors are bold lowercase ($\bm{x}$); matrices are bold uppercase ($\bm{A}$);
sets are calligraphic ($\mathcal{F}$). See also the Symbol Ledger maintained
during writing.

% =============================================================================
\section{Preliminaries}\label{sec:prelim}

\begin{definition}\label{def:lipschitz}
A function $f : \R^n \to \R$ is \emph{$L$-Lipschitz} if
\begin{equation}\label{eq:lipschitz}
\abs{f(\bm{x}) - f(\bm{y})} \le L \norm{\bm{x} - \bm{y}}
\quad \text{for all } \bm{x}, \bm{y} \in \R^n.
\end{equation}
\end{definition}

% =============================================================================
\section{Main Results}\label{sec:main}

\begin{theorem}\label{thm:main}
Let $f : \R^n \to \R$ be $L$-Lipschitz. Then \dots.
\end{theorem}

\begin{proof}
\textbf{Step 1 (Setup).}
\dots

\textbf{Step 2 (Estimate).}
By \eqref{eq:lipschitz}, \dots
\end{proof}

% =============================================================================
\section{Discussion}\label{sec:discussion}

% =============================================================================
\section*{Acknowledgements}
% Remove for double-blind submission.

% =============================================================================
\bibliographystyle{plain}
\bibliography{refs}

\end{document}
